\documentclass{article}
\usepackage{graphicx} % Required for inserting images
\usepackage{hyperref}
\usepackage{xcolor}

\title{The Meek STV algorithm, Rounding, and New Zealand rules}
\author{Andrew Conway}
\date{May 2026}

\begin{document}

\maketitle

\section{Meek in theory}

The Meek algorithm\cite{meek_nouvelle_1969,meek_1969_en,meek_1970_en} is simpler conceptually than more conventional STV algorithms
as there are fewer issues of timing etc.

Let there be $C$ candidates, $S$ seats and $T$ voters who completed valid ballots.

As with most STV algorithms, a quota $Q$ is determined as a value just high enough to
ensure that no more than $S$ candidates can exceed it if the $T$ votes are divided
between them, using the following formula, where $\epsilon$ is a small value:
\begin{equation} \label{eq:quota_original}
    Q = {{T}\over{S+1}} + \epsilon
\end{equation}

The algorithm then consists of some number of steps, each step distributing the 
votes between the candidates. If at any point a candidate gets more than $Q$, that
candidate is elected. Once $C$ candidates are elected, the count is done.

Each candidate is assigned a keep value, $0\leq k_i\leq 1$ where $i$ indicates the candidate.
Consider a ballot $b$ to have first preference for candidate $b_1$, second preference
for candidate $b_2$, etc. The vote starts off with value $1$, which gets passed to
candidate $b_1$. This candidate keeps $k_{b_1}$ of this value - it is added to that
candidate's tally, and passes on the remaining $1-k_{b_1}$ value to $b_2$.
Candidate $b_2$ then keeps $k_{b_2}$ of this $1-k_{b_1}$ value and passes the
remaining $(1-k_{b_1})(1-k_{b_2})$ value on to $b_3$. This continues until the
ballot is exhausted, and the remaining $(1-k_{b_1})(1-k_{b_2})...$ value (if any)
is counted as exhausted. That is, the value of a vote on ballot $b$ for candidate $m$
is
\begin{equation} \label{eq:vm}
    v_m = (1-k_{b_1})(1-k_{b_2})...(1-k_{b_{j-1}})k_{b_j=m}
\end{equation}
and the number of exhausted votes is
\begin{equation} \label{eq:exhaust}
    (1-k_{b_1})(1-k_{b_2})...
\end{equation}
Unlike most STV algorithms, the quota is then updated,
subtracting the number of exhausted votes $E$ from $T$:
\begin{equation} \label{eq:quota}
    Q = {{T-E}\over{S+1}} + \epsilon
\end{equation}


The keep values are determined by the following process:
\begin{itemize}
    \item All keep values start off at 1.
    \item Elected candidates (with a tally $\geq Q$) have their keep value reduced such that the candidates each end up with a tally of $Q$ after recounting with the new keep values.
    \item If an insufficient number of candidates have been elected by the prior steps,
    the candidate with the lowest tally is excluded: their keep value is set permanently
    to zero, and the tallies are recomputed.
\end{itemize}
This process continues until $S$ candidates are elected.

Determining which keep values are appropriate to make the elected candidates have a tally
of $Q$ each is not an obvious task: this involves solving a set of nonlinear
simultaneous equations. Fortunately there is a simple iterative step which
can be continued to arbitrary accuracy:
\begin{itemize}
    \item For each elected candidate $i$ with a tally $t_i\geq Q$, with keep value $k_i$ compute
    a new keep value $n_i$ via \begin{equation}\label{eq:new_keep_value}
        n_i = k_i \times {Q\over t_i}.
    \end{equation}Since $k_i>0$, $Q>0$, and $t_i\geq Q>0$, the new keep value satisfies $0<n_i\leq k_i\leq 1$.
    \item Replace the old keep values $k_i$ by the new keep values $n_i$, and recompute
    tallies. Recompute the quota as the number of exhausted votes may have increased.
    \item Note that this will reduce each elected candidate's tally by $t_i-Q\geq 0$,
    and redistribute the extra vote value $t_i-Q\geq 0$ amongst other candidates,
    including possibly other elected candidates or exhausted votes. Extra votes going
    to other elected candidates will raise their tallies above the quota; extra votes
    going to the exhausted vote tally will lower the quota; both of these will mean
    that there will be candidates with tallies higher than the quota. In this case,
    repeat the prior two steps until the values converge to some required accuracy.
\end{itemize}


Note that there are some important invariants of the Meek process that are useful
for guaranteeing convergence and some other purposes that will be discussed later:
\newcounter{invariant}
\newenvironment{invariant}[1][]{\refstepcounter{invariant}\par\medskip
   \noindent\textbf{Invariant \theinvariant. #1} \rmfamily}{\medskip}

\begin{invariant}\label{iv:nonnegative}
    \emph{Tallies can't be negative.} Nothing in equations \ref{eq:vm} or \ref{eq:quota}
    can be negative.
\end{invariant}   

\begin{invariant}\label{iv:const_sum}
    \emph{The total tallies of candidates plus exhausted votes is constant.} The sum
    over all candidates of equation \ref{eq:vm} plus exhausted value is exactly 1. Each
    ballot will always provide total value 1.
\end{invariant}   

\begin{invariant}\label{iv:keep_never_inc}
    \emph{Keep values never increase.} Keep values can only decrease by equation \ref{eq:new_keep_value} which is only applied when $t_i\geq Q$. Excluding candidates also
    only decreases keep values.
\end{invariant}   

\begin{invariant}\label{iv:cant_decrease_others}
    \emph{Reducing one candidate's keep value cannot decrease other candidate's tallies.}
    The value a ballot has for a candidate (equation \ref{eq:vm}) 
    is unaffected by the keep values of other
    candidates with a lower position on the ballot and is proportional to one minus the
    keep value for candidates higher on the ballot. This means it can increase but not
    decrease.
\end{invariant}   

\begin{invariant}\label{iv:exhausted_monotonic}
    \emph{The number of exhausted votes never decreases.} Reducing keep values can
    transfer extra vote value into the exhausted category, but cannot decrease it.
\end{invariant}   

\begin{invariant}\label{iv:quota_monotonic}
    \emph{The quota never increases.} The quota decreases as the exhausted votes increase.
    As the exhausted votes never decreases (invariant \ref{iv:exhausted_monotonic}), the quota can't increase.
\end{invariant}   

\begin{invariant}\label{iv:elected_geq_quota}
    \emph{Elected candidates never drop below quota.} Dropping a candidate's keep value
    via equation \ref{eq:new_keep_value} means that their tally will drop to $Q$,
    plus any new vote weight received from other candidates' keep value reduction.
    This will be no less than the old quota (invariant \ref{iv:cant_decrease_others}). The new quota is no higher than the
    old quota (invariant \ref{iv:quota_monotonic}), so the tally will remain higher than the quota.
\end{invariant}   

\begin{invariant}\label{iv:exhausted_not_loophole}
    \emph{The extra votes released by a reduction in quota is fewer than the exhausted votes that caused it.} When the number of exhausted votes increases by some amount
    $\delta>0$, this causes the quota to decrease by $\delta\over{(S+1)}$. Each elected
    candidate is now $\delta\over{(S+1)}$ further over quota, so there are 
    $\delta\over{(S+1)}$ more votes to distribute for each elected candidate. But
    there are no more than $S$ candidates elected, so the total amount of extra votes
    to distribute is no more than $\delta{S\over{(S+1)}}<\delta$.
\end{invariant}   

\begin{invariant}\label{iv:increase_leq_surplus}
    \emph{Without exclusions, the total increase of any tally by infinite iterations of this process cannot be greater than the total surplus (sum of the excess of the tallies of elected candidates over the quota at any point).} The process will redistribute this excess over other candidates. The
    maximum effect this can have on a particular candidate will be if all this value
    goes to that candidate. If some goes to exhausted votes this will make more available
    for redistribution, but less than it took away by being exhausted. If some goes to
    some other candidate who then goes over quota, some of that can be distributed
    back to the particular candidate, but the amount of this will be less than if
    the value had gone directly to the particular candidate.
\end{invariant}   

\begin{invariant}\label{iv:dont_elect_too_many}
    \emph{No more than $S$ candidates can reach quota.} By invariant \ref{iv:const_sum} The
    total ballots remains at $T$, so the sum of candidate's tallies is $T-E$. As none
    are negative (invariant \ref{iv:nonnegative}) this means that if $S+1$ of them were at
    quota, their sum would be $(S+1)Q$ which by equation \ref{eq:quota} is $T-E+(S+1)\epsilon>T-E$.
\end{invariant}   

\begin{invariant}\label{iv:lowest_doomed}
    \emph{If the lowest tally is smaller than the second lowest tally by more than
    the total surplus (see invariant \ref{iv:increase_leq_surplus}), one can skip
    further iterations and exclude the candidate with the lowest tally}. 
    Invariant \ref{iv:increase_leq_surplus} says that, between exclusions, 
    the lowest candidate cannot increase in tally by more than the total surplus.
    Furthermore, invariant \ref{iv:cant_decrease_others} 
    says that the other (unelected) candidates' tallies cannot decrease. So
    the lowest candidate can never catch up. This allows a huge optimization in
    the algorithm, skipping the iteration to convergence, and making the transcript
    much shorter and therefore amenable to human scrutiny.
\end{invariant}

\subsection{Convergence issues}

Convergence behaviour is generally fast and reliable \cite{meek_nouvelle_1969,meek_1969_en,meek_1970_en}, helped by the 
monotonicity guarantees above, the structure of the updates, and the $S+1$
factor in equation \ref{eq:quota}. However there is no obvious theoretical answer
I am aware of to decide what the cutoff for detecting convergence should be.

The cutoff matters, as it is conceivable that one could come up with a
pathological case where different cutoffs lead to different candidates being
elected. However a sufficiently small threshold (affecting well under a thousandth of a
vote, for instance) makes this very unlikely and furthermore so close to a tie
that its effect is hard to criticize as a failure of democracy. Of course, it is
important to specify the cutoff before the election when it is non-partisan so
that no candidate can reasonably complain to have been unfairly dealt with.

So far time complexity has been largely ignored - for millions of votes, the
running time on modern computers would be measured in seconds. The optimization
suggested in invariant \ref{iv:lowest_doomed} makes this even faster, and makes
the transcript of counts much more human readable. However, it does change the initial
keep values for the next iteration step, changing the path to convergence, and,
in the (necessary) presence of a finite cutoff for accepting sufficient convergence,
may (although almost certainly not) result in a different candidate being excluded.
Again, this is due to something that is within a tiny amount. Typically this would
be vastly under one vote), so it is not a serious issue, as long as the 
process used is specified carefully in advance to prevent partisan claims.

\subsection{Finite precision arithmetic}

The algorithm above is presented using mathematical, infinite precision, numbers.
While all computations could in principle be done with exact rational arithmetic,
a finite precision is usually easier to work with on computer. Most Australian STV
algorithms round vote tallies to integers for ease (ACT rounds to six decimal places).
Even for pedantic pompous purists (like the author) there is little to gain from
exact rational tallies given an iterative algorithm that will not necessarily converge
to rational numbers, and an associated convergence cutoff that implies finite
precision.

The big issue that rounding can bring in is that the above equations are no longer
exact, and many of the mathematical invariants may not be true, depending on how
the rounding is done. In particular, rounding may violate invariant \ref{iv:const_sum}
and possibly \ref{iv:elected_geq_quota}. These invariants no longer holding can
cause other invariants, such as invariant \ref{iv:dont_elect_too_many} to not hold, 
which is disastrous for an election system.

In Australian STV elections (which do not use the Meek algorithm), rounding is almost
universally downwards. This prevents the creation of votes above the original
$T$ and this prevents the rounding from allowing more than $S$ candidates to
exceed quota.

However, for Meek, always rounding down can be a problem, as rounding down can be 
thought of as equivalent to using a lower value of a keep value than 
one should. This would mean that \ref{iv:elected_geq_quota} does not necessarily
hold, which could cause other cascading problems.

Fortunately there is a reasonable rounding strategy for Meek.

\newcounter{roundingrule}
\newenvironment{roundingrule}[1][]{\refstepcounter{roundingrule}\par\medskip
   \noindent\textbf{Rounding rule \theroundingrule. #1} \rmfamily}{\medskip}

\begin{roundingrule} \label{rr:keep}
    \emph{When computing equation \ref{eq:new_keep_value}, round up.}
    Looking at equation \ref{eq:new_keep_value}, the ratio ${Q\over t_i}$ will never
be greater than 1 or less than zero, regardless of rounding. 
So if one rounds up\footnote{Either after doing all operations or first round up ${Q\over t_i}$ then round up the multiplication},
we retain the properties that $0\leq n_i \leq k_i \leq 1$ and $n_i \times {t_i/k_i} \geq Q$.
These properties maintain invariants \ref{iv:keep_never_inc} and \ref{iv:elected_geq_quota}.
\end{roundingrule}

\begin{roundingrule} \label{rr:vm}
    \emph{When computing equation \ref{eq:vm}, 
    round down all multiplications except the last, which should be rounded up.
    Perform multiplications left to right.}
    Looking at equation \ref{eq:vm}, we can think of this as the vote value passing
through successive keep values, some staying with that candidate, and some getting
passed on. In order to retain invariant \ref{iv:const_sum}, it is sufficient to 
ensure that one of these is rounded down, and the other is rounded up. Rounding down
the kept value is equivalent to reducing the keep value which could violate 
\ref{iv:elected_geq_quota}. Rounding up the kept value is equivalent to having a slightly
higher keep value. This appears to violate invariant \ref{iv:keep_never_inc}, however it
is not as serious as it looks. In particular, the outgoing value of the multiplication
will be no lower than the prior iteration, which maintains \ref{iv:cant_decrease_others}.\footnote{I would be happier if I had a formal proof that it could not cause any problems. But it seems much better than the alternatives.}
\end{roundingrule}

\begin{roundingrule}\label{rr:exhausted}
    \emph{When computing equation \ref{eq:exhaust} for exhausted (non-transferrable) votes,  
    $(1-k_{b_1})(1-k_{b_2})...$,
    round down each multiplication and perform multiplications left to right.}
    This is the logical extension of rounding rule \ref{rr:vm}, and maintains
    invariants \ref{iv:const_sum} and \ref{iv:exhausted_monotonic}.
\end{roundingrule}




\section{New Zealand Meek rules}

Meek's STV algorithm was greatly loved by voting geeks, but as far as I
know, was unused for government elections 
until New Zealand decided to use it for some local government
elections. Bravo New Zealand! 

The regulations\footnote{\href{https://www.legislation.govt.nz/secondary-legislation/pco-drafted/2001/145/en/latest/\#DLM57125}{https://www.legislation.govt.nz/secondary-legislation/pco-drafted/2001/145/en/latest/\#DLM57125}} are unusually well
written, and are generally a good model that I would strongly recommend to 
anyone thinking of setting up their own legislation. However there are a
few ambiguities in the timing of clauses 23, 43 and 48 that should be 
made explicit, and some serious problems with clause 10, which this document
is primarily designed to describe.

As background, keep values and tallies are both kept
to nine decimal places, and the $\epsilon$
in equation $\ref{eq:quota}$ is no more than a billionth of a vote (clauses 5 and 11). 
The convergence requirement to perform an exclusion is either a total surplus
less than 0.0001, the optimization in invariant \ref{iv:lowest_doomed} 
(clauses 13 and 14), or a catch-all oracle (clause 23). This is sufficient
precision that little would be gained by higher precision.


Clause 10 says (with some minor formatting changes):
\begin{quote}
Allocate votes from each voting document according to the following scheme:

      $v_m$ = (1 – $k_1$) × (1 – $k_2$) × (1 – $k_{m-1}$) × $k_m$

where—

$v_m$ is the votes to allocate to the mth preference candidate on the voting document

$k_1$ is the keep value for the first preference candidate on the voting document

$k_2$ is the keep value for the second preference candidate on the voting document

$k_m$ is the keep value for the mth preference candidate on the voting document

and the product of each multiplication on the right hand side of the equation must be calculated to 9 decimal digits after the point and rounded up if not exact and the multiplications must be calculated from left to right.
\end{quote}

This is exceedingly explicit, and seems unambiguous, both good properties 
of legislation. 

Comparing equation \ref{eq:vm} and the equation in the regulations,
they are the same other than the missing ellipsis. Presumably the regulations
did not mean to ignore the keep value of the third, fourth, etc. preference
(doing so would mean it would be very easy for a huge number of candidates to
go over quota - far more than the number of vacancies). I assume this is
just a typo.

More concerning, the rounding rule is explicit and unambiguous, and very
similar but different from rounding rule \ref{rr:vm}.

I believe that the actual rounding rule that is used in practice
in the New Zealand elections is rounding rule \ref{rr:vm} rather than
the rule specified in clause 10\footnote{Evidence: I can match the numbers
given in transcripts if I use rounding rule \ref{rr:vm} but not if I
use the rounding rules in clause 10.}. This note is not a 
call to New Zealand to the change way they count, but rather to explain why
I would recommend anyone writing similar legislation to use rounding rule
\ref{rr:vm} rather than copying clause 10.

\subsection{Is this rounding difference a problem?} 

\label{sec:roundingproblem}

Superficially it seems insignificant - a matter of less
than a billionth of a vote seems tiny when I earlier dismissed without justification
the vastly larger 0.0001 convergence cutoff as sufficiently accurate. 

There are two main problems:
\begin{itemize}
    \item That rounding rule violates some of the invariants, which can
    cause severe problems including potentially electing more than the desired number
    of candidates
    \item The definition in clause 2 of non-transferable votes becomes ambiguous:
    \begin{quote}
        \emph{non-transferable votes} means the votes remaining untransferred when a voting document becomes exhausted
    \end{quote}    
\end{itemize}

There are two somewhat plausible interpretations in my mind (remember I am not a lawyer)
of the above definition of non-transferable votes:
\begin{itemize}
\item \emph{Approach 1:} Extrapolate the definition of clause 10:
      use equation \ref{eq:exhaust}, with the same rounding rule as clause 10.
      This has the disadvantage that the total value of votes can easily be
      greater than 1, which violates invariant \ref{iv:const_sum} and does not
      feel like an intended property.
\item \emph{Approach 2:} Maintain invariant \ref{iv:const_sum} and insist that
      each vote has value 1, by stating that the value of a vote that is
      untransferred is 1 minus the sum of all the $v_m$ values for that vote.
      This has the disadvantage that for ballots that list all candidates - for
      which one would expect there to be no untransferred votes - this number
      will often be negative violating invariant \ref{iv:nonnegative}. 
      This does not seem to be an intended property either.
\end{itemize}

Note that if rounding rule \ref{rr:vm} were used, these two approaches would produce the
same value and violate neither invariant \ref{iv:nonnegative} nor \ref{iv:const_sum}.

Both approaches violate foundational invariants used to justify various properties
of the Meek system and cause cascading problems with other invariants.

Approach 1 is a terrible idea in practice, as it can easily lead to invariant
\ref{iv:dont_elect_too_many} to be violated, electing too many candidates.
This is caused by the value of a vote being greater than 1 - by a few billionths,
per vote. This few billionths gets multiplied by the number of votes,
making the extra votes many billionths. These extra votes are easily enough to 
exceed the $\epsilon$ in formula \ref{eq:quota} which in New Zealand is a billionth.
A practical example of this problem is given in figure \ref{fig:too_many_elected}.

Approach 2 is a bad idea in practice. While I cannot see how it could cause the
election to fail as badly as approach 1, it ends up violating at least invariants \ref{iv:nonnegative},
\ref{iv:exhausted_monotonic}, \ref{iv:quota_monotonic}, \ref{iv:elected_geq_quota},
\ref{iv:increase_leq_surplus}, and \ref{iv:lowest_doomed}. It is not clear
to me what problems this could cause.

\begin{figure}
{\footnotesize

\begin{tabular}{r|rrrrr}
Count & A & B & C & D & E \\ \hline
1 & {\color{violet} 10} & {\color{violet} 10} & 0 & 0 & 0\\
2  & {\color{violet} 6.40000002} & {\color{violet} 6.40000002} & {\color{violet} 7.19999998} & 0 & 0\\
3  & {\color{violet} 4.375} & {\color{violet} 4.375} & {\color{violet} 6.25000004} & 2.49999999 & 2.49999999\\
4  & {\color{violet} 4.0489796} & {\color{violet} 4.0489796} & {\color{violet} 4.23183676} & 3.83510204 & 3.83510204\\
5  & {\color{violet} 4.00624351} & {\color{violet} 4.00624351} & {\color{violet} 4.02872522} & 3.9793939 & 3.9793939\\
6  & {\color{violet} 4.0007934} & {\color{violet} 4.0007934} & {\color{violet} 4.0036372} & 3.99738802 & 3.99738802\\
7  & {\color{violet} 4.0001008} & {\color{violet} 4.0001008} & {\color{violet} 4.00046182} & 3.99966831 & 3.99966831\\
8 Exclude J & {\color{violet} 4.00001282} & {\color{violet} 4.00001282} & {\color{violet} 4.00005866} & 3.99995787 & 3.99995787\\
9 Exclude G & {\color{violet} 4.00000164} & {\color{violet} 4.00000164} & {\color{violet} 4.00000746} & 3.99999465 & 3.99999465\\
10 Exclude K & {\color{violet} 4.00000021} & {\color{violet} 4.00000021} & {\color{violet} 4.00000096} & 3.99999933 & 3.99999933\\
11 Exclude H & {\color{violet} 4.00000004} & {\color{violet} 4.00000004} & {\color{violet} 4.00000012} & 3.99999992 & 3.99999992\\
12 Exclude F & {\color{violet} 4.00000001} & {\color{violet} 4.00000001} & {\color{violet} 4.00000002} & 3.99999999 & 3.99999999\\
13 Exclude I & {\color{violet} 4.00000001} & {\color{violet} 4.00000001} & {\color{violet} 4.00000002} & 4 & 4\\
14  & {\color{violet} 4.00000001} & {\color{violet} 4.00000001} & {\color{violet} 4} & {\color{violet} 4.00000001} & {\color{violet} 4.00000001}\\
\hline
\end{tabular}
}
    % \includegraphics[width=0.5\linewidth]{}
    \caption{An example of an election disaster using approach 1 in section \ref{sec:roundingproblem}.
    There are 11 candidates, named A, B, C, ... K. There are 4 vacancies and 20 voters.
    The quota remains at 4.000000001 as there are no non-transferrable votes.
    Ten voters vote (in order) A,B,C,D,E,F,G,H,I,J,K. The other ten voters vote B,A,C,E,D,F,G,H,I,J,K.
    There are no non-transferable votes as all voters list all candidates.
    The table above show the tallies for candidates A to E, one per column. Candidates F to K are not
    shown as they are all zero. Tallies are marked in purple once the candidate is elected due to exceeding quota. The rounding column shows votes gained through rounding. 
    On the fourteenth iteration, candidates D and E both go over quota, causing
    five candidates to be elected for the four vacancies. 
    }
    \label{fig:too_many_elected}
\end{figure}

\section{Conclusion}

The New Zealand Meek regulations are generally excellent and I would recommend them to anyone considering STV algorithms. 

I recommend changing clause 10 to add an ellipsis in the 
formula, and change the rounding rule to:
\begin{quote}
and the product of each multiplication on the right hand side of the equation must be calculated to 9 decimal digits after the point and rounded down if not exact for all but the last, and up if not exact for the last and the multiplications must be calculated from left to right.
\end{quote}

I also recommend making timing more explicit in clauses 23, 43 and 48, or using a different tie resolution mechanism.


\bibliographystyle{plain} % We choose the "plain" reference style
\bibliography{meek} % Entries are in the refs.bib file

\end{document}
